% !TeX root = ../thuthesis-example.tex

\chapter{测试背景与目标}
\section{系统简介}
多模态方言声学智能学习平台是一个专门面向汉语方言学习与研究的综合性系统。其核心在于构建了一个内容丰富、结构严谨的“多模态方言声学知识图谱”，并基于该知识图谱开发了一系列智能算法模型，如“文脉智能体”协作框架、自适应多步检索生成模型和声学自进化优化模型等。知识图谱并非简单的词汇库，而是深度融合了语言学理论、语音声学特征和多媒体信息的认知中枢。它包含词汇、音节、声学规则、发音人、多模态资产等核心实体，以及这些实体之间的关系，例如词汇与音节的关联、音节与音频/声谱图特征的关联等。基于此知识图谱，本项目实现了“文脉智能体”框架），能针对复杂的方言问题进行链式的多智能体协作推理 ；还设计了自适应多步检索生成模型和自进化校准模型，使系统能根据学习者反馈不断优化自身性能。
\subsection{测试目的与意义}
本次测试旨在验证上述系统在真实场景下的有效性和稳定性。具体目标包括：
\begin{enumerate}
    \item \textbf{功能正确性}：验证文脉智能体是否能够准确理解并回答涉及方言声学知识的复杂问答，验证自适应生成模型是否能针对用户发音给出正确的诊断和辅导建议。
    \item \textbf{性能指标}：评估系统在自动问答和发音诊断任务上的客观性能指标，如准确率、F1值、召回率等。通过与多种基线模型对比，量化新模型的改进幅度。
    \item \textbf{个性化有效性}：考察系统根据不同学习者水平和偏好进行内容调整的能力，评估其辅导建议的完整性和实用性（通过人工评分）。
    \item \textbf{可解释性与知识覆盖}：验证系统基于知识图谱提供答案的可解释性，例如答案中引用多模态证据（音频、声谱图）的情况，以及知识图谱引用的准确性和丰富度。
    \item \textbf{自进化能力}：在持续交互过程中，检测声学自进化模型能否通过“数据生成-质量过滤-模型微调”闭环，不断提升诊断反馈的质量和准确率。
\end{enumerate}

通过以上测试，可确定整个平台是否达到了设计要求（智能化、精准化、个性化、可解释），并找出可能的不足，为后续优化提供依据。
\section{测试范围}
根据平台架构设计，本次测试主要涵盖以下四个方面。在\textbf{多模态方言声学知识图谱}部分，测试将重点关注其所包含的方言词汇、音节、声学规则、发音人及其音频、谱图等多模态信息的完整性与查询准确性，并评估其在问答与诊断任务中的支撑作用。在\textbf{文脉智能体框架}方面，将验证音韵学智能体、词汇学智能体和推理聚合智能体等多个代理（Agent）角色的分工协作效果，以及链式通信机制在传递多模态证据时的有效性。再在\textbf{自适应多步检索生成模型}的测试中，将考察其根据学习者状态动态检索知识并生成个性化答案的能力，覆盖子查询生成、子答案生成以及最终答案整合等环节。在\textbf{声学自进化生成与校准模型}部分，将评估其利用用户交互数据不断自我改进的能力，包括数据增广与模型微调机制对诊断准确率的提升效果，以及闭环自适应过程的稳定性。


\section{测试环境}
\subsection{硬件环境}
所有测试均在一台云服务器上进行，保证测试过程的独立性和可控性。该服务器配置如表\ref{tab:配置}所示。
\begin{table}[H]\centering
\caption{测试服务器硬件配置}
\label{tab:配置} % 放在caption后面
\begin{tabular}{ll}
\toprule[1pt]
\textbf{组件} & \textbf{规格} \\
\midrule[0.5pt]
处理器 & Intel Core i9-14900K \\
内存   & 32GB RAM \\
加速器 & NVIDIA Tesla 系列 GPU 1 块（用于深度学习模型推理和微调） \\
存储   & 2TB SSD（用于存储知识图谱数据库和多媒体文件） \\
\bottomrule[1pt]
\end{tabular}
\end{table}

上述配置能够满足大规模知识查询和深度模型推理的计算需求。在测试过程中，CPU主要用于知识库查询与传统算法处理（如 BM25 检索），GPU用于运行模型和语音识别、特征提取等深度学习任务。

\subsection{软件环境}
服务器操作系统为 Ubuntu 20.04 LTS 64位。主要软件和依赖如表\ref{tab:tech-components}所示。
\begin{table}[H]\centering
\caption{系统主要技术组件}
\label{tab:tech-components}
\begin{tabular}{lp{10cm}}
\toprule[1pt]
\textbf{组件} & \textbf{说明} \\
\midrule[0.5pt]
知识图谱数据库 & Neo4j 图数据库，用于存储和查询方言声学知识图谱。 \\
语言模型 & 采用讯飞星火 X 中文大语言模型作为核心模型，同时使用讯飞星火的图片理解模型对图片信息进行识别。 \\
语音处理 & 使用 Wav2Vec2-Base 预训练模型提取音频特征；使用 Montreal Forced Aligner 进行音频与文本对齐（用于验证录音对齐正确性）。 \\
检索工具 & 采用 \texttt{bm25} 库实现关键词检索 BM25 算法，\texttt{Faiss} 库用于向量相似度检索。 \\
开发框架 & Python 3.12 + PyTorch 深度学习框架。 \\
\bottomrule[1pt]
\end{tabular}
\end{table}

测试主要在离线环境下执行。本地部署了知识图谱和所需模型，测试过程中不依赖外部互联网连接，以避免网络波动影响结果。唯有在调用讯飞星火接口时，使用网络连接。整个测试环境已确保隔离稳定，避免了外部因素干扰。

\chapter{测试方法与策略}

\section{测试类型}
本次测试采用多种测试类型相结合的策略，以全面评估系统的性能与稳定性。在功能测试方面，对平台的核心功能逐项进行验证，包括知识图谱查询的正确性、智能体回答流程的完整性以及发音诊断生成的准确性。

在性能测试中，度量系统在标准数据集上的客观指标，如问答任务的准确率、召回率、F1 值，以及发音诊断任务的识别准确率和响应时间等。

在算法评估环节，重点考察各算法模块的有效性，例如比较文脉智能体框架与传统单智能体或无知识增强模型的差异，分析自适应生成模型在不同学习者状态下的输出变化，以及评估声学自进化模型在多轮交互后的性能提升。通过对比实验，将本项目模型与多种基线模型进行对照，以量化所提方法的优势。例如，将文脉智能体与直接使用 LLM（无知识库）或标准 RAG 方法进行问答表现对比，或比较自进化前后的发音诊断模型差异。通过消融实验，在移除或关闭部分模块的情况下观察性能变化，从而评估各模块对系统整体性能的贡献，例如取消音韵学智能体（合并为单 Agent）或移除多模态重排器，并比较其对问答准确率的影响。

\section{测试数据来源}
本次测试主要使用以下三类数据集，以确保评估过程既具代表性，又能够覆盖系统的核心功能与潜在边界情况。
\begin{itemize}

	\item \textbf{Dialect-QA 方言知识问答数据集}：由语言学专家基于知识图谱内容人工构建的500道问答对。问题涵盖音韵规则解释、跨方言对比等类型，答案为专家提供的标准解答，用于评估问答性能。

    \item \textbf{Pronunciation-Diag 发音诊断数据集}：收集了 200 条带有典型发音错误的学习者录音及对应标准发音。每条录音由专家标注正确的声学诊断和纠正建议，作为评估模型生成的标准答案。
    \item \textbf{合成/模拟数据}：在少量情形下，为了测试系统的异常处理能力，构造了一些边界情况的数据。例如，不存在于知识图谱的方言词条查询、极端噪声环境的录音等。这些数据用于健壮性测试，但不计入主要性能指标。
\end{itemize}

此外，知识图谱本身的数据来源于权威的方言词典、地方志及专业著作，发音音频来自方言母语者在专业录音环境下录制的语料。这些高质量数据保证了测试的可信度。

\section{测试设计思路}
测试设计遵循严谨、全面的原则，主要包含以下几个方面的考虑。

在典型场景覆盖方面，测试选取了具有代表性的使用情境。例如，在问答模块中，既设计了简单直问（单一知识点）的问题，也包含需要多跳知识推理的复杂问题；在发音诊断模块中，既测试单字发音错误的纠正，也评估连续语流中的音变检测能力。

在基线对比方面，每项测试均设置对照组，以量化改进效果。例如，在 Dialect-QA 任务上，分别使用基线模型（如 Vanilla LLM、Standard RAG、Single-Agent）与文脉智能体回答相同问题集，并比较准确率和答案质量。

在消融与单元测试方面，将多模块系统拆解为独立部分进行功能验证，如检查知识图谱查询结果是否包含预期三元组，以及智能体决策的奖励函数是否按预期工作；同时，通过消融实验分析模块交互效应，例如移除跨模态重排器对最终诊断准确率的影响。

在评估方式上，自动评估与人工评估相结合，前者提供客观指标，而对于答案解释性、辅导建议实用性等主观性较强的指标，则由三位具备语言学背景的专家进行打分，以增强评估结论的可信度。最后，针对自进化模型，设计了多轮人机交互模拟测试，每轮让模型诊断一批新录音并进行自学习，然后在独立测试集上评估性能提升，分析其随迭代次数变化的趋势，从而验证持续学习能力。

\section{测试用例设计}
表\ref{tab:test-cases} 所示为本次系统评估中选取的部分典型测试用例。这些用例覆盖了平台的两大核心功能模块——知识问答与发音诊断，任务类型包括基于声学规则的解释类问答、跨方言词汇比较、多模态证据引用等复杂推理问题，以及针对单字与多字词的发音错误检测与个性化纠正建议生成。表中列明了用例编号、任务名称、输入示例、预期结果及测试数据说明，便于实现测试过程的可复现性与结果的可解释性。

需要特别指出的是，表中仅展示了 4 个具有代表性的测试实例，旨在直观呈现测试任务类型与评估标准的多样性。实际测试过程中，我们使用了规模更大的验证数据集，总体测试执行次数超过 150 次，涵盖了不同复杂度、不同方言类型及多种噪声条件下的系统表现。

本设计确保了测试结果的代表性与统计有效性：一方面，典型用例提供了可追溯的评估样例；另一方面，大规模批量测试则保证了性能评价的稳健性和结论的普适性。
\begin{table}[H]
\centering
\caption{测试用例列表（示例）}
\label{tab:test-cases}
\begin{tabularx}{\textwidth}{p{1.8cm} p{2.2cm} p{3cm} X X}
\toprule
\textbf{用例编号} & \textbf{名称} & \textbf{输入} & \textbf{预期结果} & \textbf{测试数据说明} \\
\midrule
QA-001 & 苏州话入声弱化现象解释 & 用户提问：“苏州话中‘袜’一词的读音为何会发生入声弱化？” & 系统解释韵尾弱化，并提供苏州话“袜”标准发音音频链接作为佐证。 & 涉及声学规则（入声弱化），需结合知识图谱与音频档案。 \\
\midrule
QA-002 & 吴语与闽南语对应词比较 & 用户提问：“分析‘银行’在苏州话中的声母音变及与闽南语的差异。” & 指出声母变化并给出 IPA 符号；提供闽南语对应词发音及差异说明，并附音频证据。 & 涉及多跳推理：分别检索苏州话和闽南语的声学数据，多模态资源辅助比较。 \\
\midrule
PD-001 & 入声韵尾缺失的发音诊断 & 学习者发“入(ru4)”字，韵尾[-k]缺失。 & 系统诊断韵尾缺失并给出收声练习建议，附标准音及声谱图对比。 & 针对入声韵尾缺失错误，需比对标准音与用户音定位问题。 \\
\midrule
PD-002 & 连续变调错误的诊断 & 学习者读上海话“老早”未变调。 & 系统诊断缺少连读变调并说明规则，附正确发音示例。 & 检验变调识别能力，需结合知识图谱规则与音档例句。 \\
\bottomrule
\end{tabularx}
\end{table}
\chapter{测试过程与执行}
\section{执行步骤}
在测试执行阶段，整个过程严格依照预先制定的用例计划逐一开展。在环境准备方面，需要启动知识图谱数据库并导入最新的多模态数据，同时加载或部署各类智能模型，包括大语言模型（LLM）、检索器以及诊断生成模型，确保所用模型均采用训练完成的参数，尤其是自适应模型在微调后的最新版本。 
 
在执行用例时，对于问答类任务（QA-001 与 QA-002），通过系统问答接口提交问题文本，并观察文脉智能体框架的内部运行过程。该框架会首先由调度智能体解析问题，将任务分别分配给音韵学 Agent 和词汇学 Agent，在知识图谱（$D_{\mathrm{kg}}$）与方言音档（$D_{\mathrm{archive}}$）中检索所需信息。各 Agent 检索完成后会生成包含文本摘要与多模态证据 URI 的通信单元，并依次传递给下一个 Agent；最终，由管理 Agent 汇总全部证据，完成多模态链式推理并生成答案。对于发音诊断类任务（PD-001 与 PD-002），通过系统的发音诊断接口上传用户录音，系统首先利用 Wav2Vec2 模型提取声学特征向量，然后在知识库中检索匹配的标准发音音频及其特征。该检索过程先通过 Faiss 向量相似度检索生成候选，再由跨模态重排器进行精细匹配以确定最佳标准音。接着，自适应生成模型以用户音频与标准音为输入，生成包含诊断结果与针对性的辅导建议的文本输出。  

在结果记录阶段，测试人员会保存系统返回的答案文本及多媒体证据，例如音频链接与声谱图截图等。同时，如系统提供日志或可观察的内部状态（如 Agent 检索到的知识片段、奖励函数打分情况），也会一并记录，用于后续分析系统行为与预期的一致性。若在执行过程中出现系统未返回结果或发生错误（例如因输入格式问题导致解析失败），则会记录问题并根据用例要求调整输入后重新执行。例如，对于涉及多语言的 QA-002，用例可能在首次执行时得到不完整回答，此时会适当细化提示并再次测试，以判断问题源于系统逻辑还是用例本身的设计不当。

整个执行过程中，我们严格按照预定步骤，确保每个用例至少执行一次完整流程，必要时重复验证其一致性。
\section{实际执行情况记录}
大部分测试用例的执行结果均符合预期，系统功能表现稳定。

对于问答用例，文脉智能体框架能够顺畅地完成链式推理。如在用例QA-001中，音韵学Agent成功从知识图谱检索到“入声弱化”规则描述，词汇Agent找到“袜”在苏州话中的发音音频URI，两者信息在最终答案中被有效融合，系统答案准确解释了入声弱化成因并附上音频证据。内部日志显示，每个Agent的检索决策均获得了较高的相关度奖励$\mathcal{R}_{rel}$，且避免了重复检索。这验证了奖励函数引导Agent检索行为的有效性。

\begin{figure}[H]\centering
\includegraphics[width=0.9\textwidth]{QA1.png}
\caption{QA-001部分日志}
\end{figure}

在用例QA-002中，系统生成的答案包含了苏州话“银行”的IPA和音变说明，以及闽南语对应词的对比，和预期相符。值得一提的是，答案中特意引用了知识图谱中“尖团合流”音变规则来解释苏州话声母变化，并在结尾提供了苏州话“银行”一词的标准读音音频链接。这说明智能体不仅检索到表面答案，还综合调用了背景声学知识，增强了解释性。实际执行中，多模态证据链以URI形式嵌入通信单元在Agent间传递，最终由管理Agent转化为答案中适当的文字和引用。这一链式通信机制运行如预期，没有出现证据丢失或遗漏情况。

对于发音诊断用例，系统同样表现良好。用例PD-001中，系统准确检测出用户“入”字录音缺少喉塞韵尾，并给出了纠正建议和标准音频。观察日志可见，跨模态重排器将用户音频与知识库中多个候选标准音进行比对后，成功将正确的“入”字标准发音排在首位。生成的诊断文本几乎与专家标注的标准答案一致，证明自适应生成模型有效地学到了专业诊断表述。在用例PD-002中，系统检测连读变调错误稍有延迟（因为需分析完整语音片段约2秒），但最终输出内容完整正确，指出了问题并教授了正确声调变化，附带的参考音频播放验证了建议的准确性。

\begin{figure}[H]\centering
 \includegraphics[width=0.9\textwidth]{用例3日志.png}
\caption{PD-002部分日志}
\end{figure}
在整个执行过程中，未发生系统崩溃或长时间无响应的现象。测试执行过程顺利，大部分用例一次通过。少数需要关注的是，个别复杂问题（涉及跨多个知识关联）系统第一次回答略有遗漏，通过提示重新提问后得到完善答案。这提示我们未来可考虑增强智能体在复杂查询计划上的完备性。总体而言，实际执行结果证明系统功能达到了设计预期。



\section{测试结果与分析}

经过严格的测试执行，我们对系统在Dialect-QA问答任务和Pronunciation-Diag发音诊断任务上的表现进行了量化评估，并对比了多个基线模型。同时，通过消融实验分析了各关键模块的重要性。以下将分别汇报自动评估指标、人工评估结果以及相关数据可视化分析。
\subsection{自动评估指标}
\textbf{问答任务性能：}在Dialect-QA数据集上，本项目模型得了显著优于基线的成绩。如表7.1所示，Dialectus在准确率上达到了\textbf{55.3\%}，较直接使用LLM的15.4\%提高了近40个百分点；F1值达到\textbf{76.8\%}，也远高于其他对照模型 。尤其值得注意的是，传统检索增强模型(Standard RAG)的F1仅为65.7\%，说明简单拼接知识的方案未能充分利用知识，而Dialectus通过链式推理充分发挥了知识图谱的作用 。此外，Dialectus在长答案上的ROUGE-L指标为62.4，显示其给出的解释性答案与标准答案有高度重合，内容全面详实 。

\begin{table}[h]\centering
\caption{Dialect-QA问答任务自动评估结果}
\begin{tabular}{lccc}
\toprule[1pt]
模型 & EM(\%) & F1(\%) & ROUGE-L \\
\midrule[0.5pt]
Vanilla LLM (零样本) & 15.4 & 23.1 & 20.5 \\
Standard RAG & 42.8 & 65.7 & 51.3 \\
Single-Agent (简化智能体) & 48.1 & 60.2 & 48.9 \\
\textbf{Dialectus (本项目)} & \textbf{55.3} & \textbf{76.8} & \textbf{62.4} \\
\bottomrule[1pt]
\end{tabular}
\end{table}

从上表可以看出，简化的Single-Agent模型（将所有知识一次性提供LLM）比Standard RAG略有提升，但与Dialectus仍有明显差距。这说明仅提供知识是不够的，关键在于如何使用知识。Dialectus的多Agent分治协作有效地解构了问题，并引入了多模态证据链，使模型理解更深入全面 。

\textbf{发音诊断任务性能：}在Pronunciation-Diag数据集上，自进化后的诊断模型表现卓越。表7.2显示，模型微调前（Generator$_{\theta_0}$）的诊断准确率F1为68.3，而经过多轮交互自学习后的模型(Generator+Self-Evolve) F1提升至\textbf{89.6}，提高了21.3个百分点 。这证明自进化闭环极大增强了模型对发音错误的识别和纠正能力。此外，我们对模型诊断的细分类准确率做了统计，例如声母错误识别F1从64提升到87，声调错误识别F1从72提升到91，均有显著改善。模型在保持高准确率的同时，生成反馈的平均长度和丰富度也增加了约30\%，说明其辅导内容更加详尽。

\begin{table}[H]\centering
\caption{Pronunciation-Diag 发音诊断任务自动评估结果}
\begin{tabular}{lccc}
\toprule[1pt]
模型 & 总体F1(\%) & 声母错误F1(\%) & 声调错误F1(\%) \\
\midrule[0.5pt]
Generator$_{\theta_0}$（微调前） & 68.3 & 64 & 72 \\
\textbf{Generator+Self-Evolve（本项目）} & \textbf{89.6} & \textbf{87} & \textbf{91} \\
\bottomrule[1pt]
\end{tabular}
\end{table}

从结果看，自适应微调使模型对专家标注的“黄金标准”学习到了更加贴近的诊断模式，特别是在一些微妙的音变错误上，进化后的模型能够给出几乎与专家描述相同的解释。这也体现了RAFT增强微调方法的威力：模型学会了如何“利用参考资料进行思考”，将知识图谱提供的标准答案内化为了自身的推理能力。

\subsection{人工评估结果}
为了评估答案的质量和辅导建议的实用性，我们邀请了三位语言学专业人士对模型输出进行人工评分，每项5分制。评价维度包括问答的正确性、完整性以及诊断辅导的有效性。

在Dialect-QA问答任务上，Dialectus生成的答案平均正确性评分为4.88/5，完整性为4.81/5，远高于基线模型（Vanilla LLM仅约3.1和2.9，Single-Agent约4.3和4.2）  。专家评委一致认为，Dialectus的回答几乎没有明显错误，且往往提供额外的背景解释和例证，使答案更加全面。如在“银行”声母音变的问题上，Dialectus不仅回答出声母变化是什么，还进一步解释了该变化的历史声学原因，并引用知识图谱规则，体现出很强的专业性和可信度。

在发音诊断任务上，进化后的模型辅导有效性评分为4.72/5，相比初始模型的3.55有大幅提升 。评委反馈指出，新模型的反馈更具针对性和可操作性。例如，用例PD-001中，模型不仅指出错误，还给出了具体练习方法（如何发喉塞音），并提供了标准音频供模仿，具有良好的教学效果。相较之下，初始模型有时只给出笼统的纠正建议，缺乏细节。由此可见，自进化模型通过不断学习，已经能产出接近专家水平的个性化辅导内容。

\subsection{数据表格与图表分析}
通过对比实验和消融实验，我们进一步分析各模块的贡献和系统性能瓶颈。

\textbf{基线对比分析：}从表\ref{tab:main_results_qa}和表\ref{tab:main_results_diag}的数据可以明显看出，本项目的方法在两项任务上都取得了最好的效果。特别是在复杂的Dialect-QA任务上，文脉智能体(Dialectus)的准确率和综合得分均拔得头筹。这说明引入结构化知识和多模态推理对于深度问答是非常有效的。另一方面，Standard RAG虽然结合了知识库，但其性能甚至低于Single-Agent，这可能是因为RAG简单拼接文本导致LLM产生“信息过载”或无法有效利用知识。而Single-Agent有所改进，意味着在有知识辅助的情况下LLM表现提升，但仍不如我们的链式多Agent方案。
\begin{table}[H]
    \centering
    \caption{在Dialect-QA数据集上的模型性能对比}
    \label{tab:main_results_qa}
    \begin{tabular}{@{}l|cc|c|cc@{}}
        \toprule
        \multirow{2}{*}{\textbf{模型}} & \multicolumn{2}{c|}{\textbf{自动评估}} & \multicolumn{3}{c}{\textbf{人工评估 (1-5分)}} \\
        & EM & F1 Score & ROUGE-L & 正确性 & 完整性 \\
        \midrule
        Vanilla LLM & 15.4 & 23.1 & 20.5 & 3.12 & 2.89 \\
        Standard RAG & 42.8 & 55.3 & 48.9 & 4.15 & 4.05 \\
        Single-Agent & 48.1 & 60.2 & 51.3 & 4.33 & 4.21 \\
        \textbf{Dialectus (本项目)} & \textbf{65.7} & \textbf{76.8} & \textbf{62.4} & \textbf{4.88} & \textbf{4.81} \\
        \bottomrule
    \end{tabular}
\end{table}

\begin{table}[H]
    \centering
    \caption{在Pronunciation-Diag数据集上的模型性能对比}
    \label{tab:main_results_diag}
    \begin{tabular}{@{}l|c|c@{}}
        \toprule
        \textbf{模型} & \textbf{Diagnosis Acc. (F1)} & \textbf{辅导有效性 (1-5分)} \\
        \midrule
        Generator($\theta_0$) & 68.3 & 3.55 \\
        \textbf{Generator+Self-Evolve (本项目)} & \textbf{89.6} & \textbf{4.72} \\
        \bottomrule
    \end{tabular}
\end{table}


\textbf{消融实验分析：}本研究针对 Dialectus 框架与自进化模型分别进行了关键组件的消融测试，结果如图\ref{fig:消融}所示。实验表明，移除部分组件会导致性能出现显著下降。首先，当不使用音韵学智能体，而是将其与词汇 Agent 合并为单一 Agent 处理所有信息时，Dialect-QA 的 F1 值下降了 16.6 个百分点，这是所有设置中降幅最大的一项。这一结果表明，专门的音韵学 Agent 在处理声学知识方面发挥着不可替代的作用，其能够有针对性地检索和分析复杂的声学规则，从而显著提升答案的准确性。其次，在禁用链式通信机制的情况下（即各 Agent 并行独立回答并在最后简单合并结果），F1 值下降了约 12.3 个百分点。这说明逐步传递信息、让后续 Agent 基于前一步的总结继续推理，比并行处理更能生成连贯且深入的答案。链式结构的引入保证了每一步推理都有上下文支撑，从而减少了答案的片面性。

在发音诊断任务中，移除多模态重排器（仅使用第一阶段混合检索结果而不进行交互精排）会使诊断准确率下降 7.4 个百分点。虽然系统仍能找到较为相关的候选标准音，但缺少精细匹配容易导致选错参考音频，进而影响诊断结果的正确性。这表明跨模态重排能够有效利用用户音频与候选音频之间的细粒度交互信息，提高检索精度。同样地，若取消教师模型过滤（即不对生成数据进行质量筛选而直接用于微调），诊断任务的 F1 值会下降约 9.1 个百分点，这意味着未经过滤的训练数据中可能混入错误反馈，从而使模型学习到不良模式。教师模型在此过程中确保了只有高质量样本才能用于自学习，是提升模型诊断能力的重要保障。

最后，当关闭奖励函数引导（使 Agent 在检索时随机或贪心选择信息，而不考虑信息增益）时，Dialect-QA 的性能下降了 5.8 个百分点。尽管核心 Agent 依然存在，但检索行为失去了优化策略，会出现重复检索或偏离主题的情况，从而降低了检索效率与答案的准确率。奖励函数在关联度、新颖度和多模态价值等维度的引导作用，确保了 Agent 的每一步操作都能带来较高的信息收益。
\begin{figure}[H]\centering
 \includegraphics[width=0.9\textwidth]{消融.pdf}
\caption{关键模块消融对性能的影响：移除各组件后Dialect-QA F1下降情况（性能下降绝对值越大表示该组件贡献越大）。}
\label{fig:消融}
\end{figure}

为了进一步验证声学自进化模型的有效性和收敛性，我们追踪了模型在多轮自进化过程中的性能变化，结果如图\ref{fig:pronunciation_diag_convergence}所示。

\begin{figure}[H]
    \centering
    \includegraphics[width=\textwidth]{f1.pdf}
    \caption{基于表3.4数据的声学自进化模型收敛曲线。图中展示了Generator+Self-Evolve模型从初始Generator$_{\theta_0}$的68.3\%提升到最终89.6\%的F1分数，实现了21.3个百分点的性能增益。浅蓝色区域表示95\%置信区间。}
    \label{fig:pronunciation_diag_convergence}
\end{figure}

图中蓝色实线代表Generator+Self-Evolve模型的F1分数，随着自进化轮次的增加呈现出显著的单调递增趋势。这表明，本文提出的"交互-生成-过滤-学习"自进化闭环是有效的，模型能够持续地从用户反馈中汲取知识，实现性能的自我提升。自进化模型从初始Generator$_{\theta_0}$的68.3\%起步，在第1轮自进化后便显著超越了初始性能，并最终收敛于89.6\%的水平，实现了21.3个百分点的显著提升。这定量地证明了自进化机制相较于静态模型的优越性。

图中的浅蓝色阴影区域代表了95\%置信区间。该区间在训练初期相对较宽，反映了模型学习初期的不确定性；但随着自进化轮次的增加，置信区间逐渐变窄，表明模型的性能表现愈发稳定。这证明了模型的收敛过程是稳健可靠的，而非偶然的性能波动。绿色填充区域直观地展示了自进化带来的性能增益，清晰地量化了本方法的有效性。

综上所述，该图从趋势、终值和稳定性三个维度，强有力地证明了本文所提出的声学自进化方案是一个有效、优越且稳健的持续学习框架，能够实现从68.3\%到89.6\%的显著性能提升。

从以上分析可知，音韵学Agent、链式通信在问答任务中属于\textbf{关键组件}，多模态重排器和教师过滤在各自任务中为\textbf{重要组件}，而奖励函数虽作用相对稍小但也不可忽视。这样的实验定量证明了本项目设计中各模块的有效性，也解释了系统性能远超基线的原因：每个模块都在发挥作用，共同造就了整体性能的跃升。

\textbf{案例分析：}为了更直观地感受系统优势，我们选取问答任务中的典型案例进行对比（来自Dialect-QA数据集）。问题是：“苏州话‘银行’一词的声母，相比普通话发生了怎样的音变？请结合其IPA进行解释。”
传统Baseline（Standard RAG）给出的答案：“苏州话‘银行’的标准音是[ɲәŋ hɑ̃]。普通话是[yínháng]。声母有区别。”——可以看出仅点出了存在区别但未详述。
而文脉智能体框架的回答：\emph{“苏州话‘银行’[ɲәŋ hɑ̃]的声母相较普通话[yín háng]发生了显著的音变。具体而言：1）‘银’字：普通话声母为舌尖鼻音[n]（拼音“n”），而在苏州话中演变为软腭鼻音[ŋ]，这是吴语中尖音变为团音的结果。2）‘行’字：普通话声母为声门擦音[h]，苏州话中由于中古匣母读音的保留，声母为软腭鼻音[ŋ]，且受后续鼻化元音影响，[ŋ]在口语中有弱化趋势。一些苏州话口音中该字甚至读作零声母。您可以收听标准读音音频以体会这一特点。”}
可以看到，文脉智能体框架不仅指出了具体“n→ŋ”的变化，还提及了尖团合流等声学现象作为背景，解释了为什么会这样变化。同时提供的标准发音音频让用户亲耳验证差异。专家评委认为，这种回答正确、完整且具有教学意义，远胜于简单的答案。该案例体现了知识图谱和多Agent协同带来的丰富性，以及多模态输出增强了答案的可信度和实用性。

\section{结论}
本次测试充分验证了多模态方言声学智能学习平台在功能和性能上的优越性。基于测试结果和分析，我们得出以下结论：
\begin{enumerate}
    \item \textbf{系统满足设计目标}：平台能够准确、深入地解答复杂的方言知识问答，自动评估指标和人工评分均表明其回答正确率高、内容完整；在发音诊断与辅导上，系统提供的反馈专业有效，达到接近人工专家的水平。这说明系统实现了智能化、精准化和个性化的设计初衷。
    \item \textbf{知识图谱赋能显著}：测试表明，引入多模态方言声学知识图谱极大提升了问答的准确性和可解释性。系统回答能够引用知识图谱中的规则和例证，使用户获得验证过的知识来源。对比实验也证明，没有知识图谱支撑的模型表现大打折扣，知识图谱是整个系统智能的\textbf{基石}。
    \item \textbf{多 Agent 链式推理有效}：文脉智能体框架通过划分角色和分步推理，成功应对了长上下文和复杂查询。测试结果显示，它较单一 Agent 方案有明显性能优势，且答案质量更高。这验证了链式通信、小窗口处理长上下文的可行性，并展示了在专业领域应用 LLM 的新范式。
    \item \textbf{自适应个性教学可行}：系统能够根据学习者掌握情况调整输出，通过动态子查询和内容生成，实现了“因材施教”。在测试用例中，不同难度的问题系统给出的解答详略得当；对不同发音错误的用户，系统提供的纠正方案各有侧重。这表明用户状态建模和自适应内容生成策略取得了预期效果。
    \item \textbf{持续自进化提升性能}：声学自进化模型通过交互数据不断优化，显著提高了发音诊断准确率和反馈质量。测试中每轮微调都带来性能上升，证明了封闭式人机教学循环的价值。随着更多用户数据累积，预计模型性能还将进一步提升。
    \item \textbf{模块设计合理}：消融实验的结果支持了各模块设计的必要性。特别是音韵学 Agent、链式通信、多模态重排器和教师过滤等，对系统性能影响明显。这说明项目在知识获取、信息过滤、多模态融合等方面的设计决策是成功的，各部分协同运作，缺一不可。
\end{enumerate}

本测试得出结论：\textbf{多模态方言声学智能学习平台通过严谨的架构和算法设计，实现了优秀的问答和辅导性能，满足预期需求}。系统在测试范围内运行稳定、输出可靠，为实际部署应用奠定了基础。当然，测试也暴露了一些改进空间，例如对极端复杂问题的处理效率、对更多方言的扩展等。在未来的工作中，我们将根据测试发现持续优化，进一步提高系统的健壮性和泛化能力。
